%! Author = jacopo
%! Date = 2/25/26
\section{Introduction}
Machine Learning advancements can be applied across a wide range of domains due to their ability to model complex data.
This has also influenced the field of affective computing where electroencephalography (EEG) is used to infer emotional and cognitive states.
In recent years many models have been developed for EEG representation learning.
These general purpose models, so-called \textit{foundation models}, in the field tend to be limited by various reasons, with the low cardinality of EEG centric datasets, being one of them.
This is especially noticeable once one considers the size of corpora of most text-based state-of-the-art (SOTA) models.

While multimodal foundation models have shown success in vision and language, their application to EEG representation learning remains limited.
With this problem in mind the goal of the project is to leverage cross-modal knowledge to integrate information into EEG representations.
Our approach is limited to the affective aspect of neurological data processing but the core idea can be applied to any field.

EEGAVI transfers information from multiple different streams into the EEG aiming to enhance the representational capacity.
It employs a principled integration of existing techniques (cross-attention + contrastive loss + knowledge distillation) applied to EEG under data scarcity.
The model leverages multiple backbones for the elaboration of the input data sources, for EEG in particular CBraMod is employed.

Beyond the need for a foundation model the research question is mostly concerned about understanding the impacts of \textit{knowledge distillation} in these constraints.
In particular to see if distillation in a multimodal context where sources and target do not belong to the same modality is of benefit to the training pipeline.
To answer this question we pretrain an instance of EEGAVI and analyze its performance using CBraMod as a reference point.
These analyses are first performed on cross-retrieval metrics on the held-out data and later in various downstream tasks.
